\documentclass[a4paper]{article}
\usepackage[utf8]{inputenc}
\usepackage[russian]{babel}
\usepackage[margin=1in]{geometry}
\usepackage{float}
\usepackage{graphicx}
\usepackage{setspace}
\title{Лабораторная работа}
\author{Шулайкин Д.А.}
\begin{document}
\onehalfspacing
\thispagestyle{empty}
\begin{center}
Министерство образования и науки Российской Федерации
\vspace{10pt}

Федеральное государственное бюджетное образовательное учереждение высшего образования учреждения высшего профессионального образования "Ивановский государственный энергетический университет имени В.И. Ленина"
\vspace{40pt}

Кафедра програмного обеспечения компьютерных систем
\vspace{40pt}

\textbf{Отчет по лабораторной работе №3}

Дисциплина: Операционные системы

Тема: «РАБОТА В СРЕДЕ КОМАНДНОЙ ОБОЛОЧКИ MICROSOFT POWERSHELL»

Вариант 12
\end{center}

\vspace{330pt}
\begin{flushright}
\textbf{Выполнил:}

Студент гр. 2-42В, Шулайкин Д.А.

\textbf{Проверил:}

Гадалов А.Б.
\end{flushright}
\vspace{40pt}
\begin{center}
Иваново 2018
\end{center}
\pagebreak

\section{Задание}

POWERSHELL + WINFORMS
\begin{enumerate}
\end{enumerate}
\section{Выполнение работы}

Запустите виртуальную машину с предустановленной ОС MS Windows 7, подготовленную при выполнении лабораторной работы №2.
\begin{figure}[H]
    \centering
    \includegraphics[width=\linewidth]{1.png}
    \caption{start VM}
\end{figure}

Запустите командную оболочку PowerShell, познакомьтесь со справочной системой.
\begin{figure}[H]
    \centering
    \includegraphics[width=\linewidth]{2.png}
    \caption{get-help}
\end{figure}

\begin{figure}[H]
    \centering
    \includegraphics[width=\linewidth]{3.png}
    \caption{get-command}
\end{figure}

В домашней директории пользователя привычными средствами (например, через Проводник) создайте каталог Lab3. В этом каталоге (теми же средствами) создайте подкаталоги с именами изучаемых в этом семестре дисциплин (не менее 4, латиницей), а в каждом таком каталоге – подкаталоги с произвольными именами, например, Lectures, Labs, Exam (не менее 3). В каталог OS (Операционные системы) в один из подкаталогов (например, Labs) поместите не менее двух файлов с произвольным содержимым, например, скопируйте свои отчеты по лабораторным работам №1, №2.

\begin{figure}[H]
    \centering
    \includegraphics[width=\linewidth]{5.png}
    \caption{dmitrii's home directory}
\end{figure}

\begin{figure}[H]
    \centering
    \includegraphics[width=\linewidth]{6.png}
    \caption{\~/Lab3/}
\end{figure}


\begin{figure}[H]
    \centering
    \includegraphics[width=\linewidth]{7.png}
    \caption{\~/Lab3/os/labs}
\end{figure}

Отобразите на экране (средствами PowerShell) содержимое домашней директории пользователя (рекурсивно).

\begin{figure}[H]
    \centering
    \includegraphics[width=\linewidth]{4.png}
    \caption{\~/Lab3/os/labs/}
\end{figure}


\begin{figure}[H]
    \centering
    \includegraphics[width=\linewidth]{8.png}
    \caption{\~/Lab3/os/labs/}
\end{figure}

Выведите на экран информацию обо всех файлах, находящихся в каталоге /OS/Labs и удовлетворяющих некоторой маске.

\begin{figure}[H]
    \centering
    \includegraphics[width=\linewidth]{9.png}
    \caption{ls -R ./Lab3/os/labs/*2*}
\end{figure}

Получите список дисков PowerShell, доступных в текущем сеансе работы.

\begin{figure}[H]
    \centering
    \includegraphics[width=\linewidth]{10.png}
    \caption{get-psdrive}
\end{figure}

Продемонстрируйте навигацию между созданными подкаталогами домашней директории пользователя с использованием командлета Set-Location, используйте абсолютные и относительные пути.

\begin{figure}[H]
    \centering
    \includegraphics[width=\linewidth]{11.png}
    \caption{cd relative paths}
\end{figure}

\begin{figure}[H]
    \centering
    \includegraphics[width=\linewidth]{12.png}
    \caption{cd abs paths}
\end{figure}

В домашней директории пользователя средствами PowerShell создайте новый диск с именем myLabs\_OS, ассоциированный с каталогом, в котором хранятся отчеты по лабораторным работам по дисциплине «Операционные системы».

\begin{figure}[H]
    \centering
    \includegraphics[width=\linewidth]{13.png}
    \caption{new-psdrive -name myOsLabs -psprovider FileSystem -root ./Lab3/os/labs/}
\end{figure}

Посмотрите содержимое каталога с отчетами по лабораторным работам, используя полный путь к каталогу, и используя имя диска, только что созданного и ассоциированного с этим каталогом.

\begin{figure}[H]
    \centering
    \includegraphics[width=\linewidth]{14.png}
    \caption{ls ./Lab3/os/labs/}
\end{figure}

\begin{figure}[H]
    \centering
    \includegraphics[width=\linewidth]{15.png}
    \caption{ls myOsLabs:/  }
\end{figure}

Сохраните несколько скриншотов в каталоге с отчетами по лабораторным работам. Выведите список всех файлов с расширением png (или jpg), сохраненных в каталоге с отчетами по лабораторным работам. Продемонстрируйте использование параметра –Exclude.


\begin{figure}[H]
    \centering
    \includegraphics[width=\linewidth]{16.png}
    \caption{ls ./Lab3/os/labs/ -exclude *png,*tex }
\end{figure}

Выведите список только подкаталогов директории ./OS, отсортированный по времени создания. Определите размер каталога ./OS.


\begin{figure}[H]
    \centering
    \includegraphics[width=\linewidth]{18-2.png}
    \caption{ls ./Lab3/os/ -R | measure-object -property length -sum }
\end{figure}

\begin{figure}[H]
    \centering
    \includegraphics[width=\linewidth]{18-3.png}
    \caption{ sort subdirectories by creation time }
\end{figure}

Средствами PowerShell создайте новый каталог test\_dshulaikin в каталоге ./OS. В этом каталоге создайте файл, в котором сохраните время выполнения этого пункта задания (например, test\_dshulaikin). Прочитайте содержимое созданного файла. Добавьте в конец файла сообщение «Тест успешно пройден» («Test is completed successfully») и новое время выполнения этого пункта задания.

\begin{figure}[H]
    \centering
    \includegraphics[width=\linewidth]{20.png}
    \caption{ make directory }
\end{figure}


\begin{figure}[H]
    \centering
    \includegraphics[width=\linewidth]{24.png}
    \caption{ write to file }
\end{figure}

Создайте новый каталог newTest в каталоге другой дисциплины. Скопируйте в этот каталог файл test\_dshulaikin. Отобразите содержимое нового каталога (рекурсивно), чтобы убедиться, что файл скопирован.

\begin{figure}[H]
    \centering
    \includegraphics[width=\linewidth]{25.png}
    \caption{mkdir}
\end{figure}


\begin{figure}[H]
    \centering
    \includegraphics[width=\linewidth]{26.png}
    \caption{cp, ls}
\end{figure}

Переименуйте каталог.

\begin{figure}[H]
    \centering
    \includegraphics[width=\linewidth]{27.png}
    \caption{mv}
\end{figure}

Удалите файл test\_file из каталога ./OS.

\begin{figure}[H]
    \centering
    \includegraphics[width=\linewidth]{29.png}
    \caption{rm ./Lab3/os/test$\_$file.txt}
\end{figure}


Подготовьте отчет о выполненной работе.
\subsection{Подготовьте отчет о выполненной работ}


\section{Выводы по работе}

PowerShell

\end{document}
